\section{Characteristic transfer and smooth linear sections}
\label{sec:transfer}

Proposition~\ref{prop:gradient-codimension} bounds an affine gradient
locus in characteristic two.  We now transfer that bound to complex
parameters, then choose a linear section on which the coefficient
polynomial defines a smooth projective hypersurface.  The final proof
will combine the resulting affine permanent projection with
Theorem~\ref{thm:smooth-border} to obtain the cubic lower bound.

\subsection{Transferring the gradient bound}

The proper family used for specialization is the projective zero scheme
of the partial derivatives.  In characteristic two this scheme need not
be the singular scheme of the hypersurface itself.  After passing to
characteristic zero, Euler's identity will identify the two.

\begin{lemma}[Specialization of a gradient locus]
\label{lem:characteristic-transfer}
Let
\[
 F(\theta,x)\in
 \mathbb Z[\theta_1,\ldots,\theta_a][x_1,\ldots,x_v]
\]
be homogeneous of degree \(r\ge2\) in \(x\), where \(v\ge2\), and let
\(1\le\delta\le v\) be an integer.  Suppose some
\(\bar\theta\in\overline{\mathbb F}_2^{\,a}\) satisfies
\[
 \dim V_{\mathbb A^v_{\overline{\mathbb F}_2}}
       \bigl(\partial_{x_1}F_{\bar\theta},\ldots,
             \partial_{x_v}F_{\bar\theta}\bigr)
 \le v-\delta.
\]
Then there is a nonempty Zariski open subset
\(U\subset\mathbb A^a_{\C}\) such that, for every \(\theta\in U\),
the polynomial \(F_\theta\) is nonzero and
\begin{equation}
 \dim\operatorname{Sing}V_{\PP^{v-1}_{\C}}(F_\theta)
 \le v-1-\delta.
 \label{eq:transfer-singular-bound}
\end{equation}
We use dimension \(-1\) for the empty projective scheme.
\end{lemma}

\begin{proof}
Put \(R=\mathbb Z[\theta_1,\ldots,\theta_a]\) and
\(S=\operatorname{Spec}R\).  The homogeneous ideal
\[
 I=(\partial_{x_1}F,\ldots,\partial_{x_v}F)\subset R[x_1,\ldots,x_v]
\]
defines a closed subscheme
\begin{equation}
 \mathcal Z=\operatorname{Proj}(R[x]/I)\subset\PP^{v-1}_S.
 \label{eq:transfer-projective-family}
\end{equation}
Its structural morphism \(\pi:\mathcal Z\to S\) is projective, hence proper.
Differentiation commutes with specialization, and the closed projective
subscheme in \eqref{eq:transfer-projective-family} commutes with base
change.  Thus its geometric fibers are exactly the projective
gradient-zero schemes of the specialized polynomials.

The derivatives have positive homogeneous degree \(r-1\).
At \(\bar\theta\), their affine zero set is a cone.  If its
projectivization is nonempty, its dimension is one less than that of the
cone: on each projective coordinate chart the corresponding punctured
cone is a product with \(\mathbb G_m\).  If the projectivization is
empty, the same upper bound holds with our convention.  Hence
\[
 \dim\mathcal Z_{\bar\theta}\le v-1-\delta.
\]
Let \(s\in S\) be the image of the geometric point \(\bar\theta\).
The dimension of a scheme of finite type over a field is unchanged by
extending that field, so the same bound holds for \(\mathcal Z_s\).
If \(v-1-\delta\ge0\), upper semicontinuity of fiber dimension for the
proper morphism \(\pi\) gives an open neighborhood \(U_S\) of \(s\)
on which this bound holds \cite[Tag~0D4I]{StacksProject}.
If \(v-1-\delta=-1\), the fiber \(\mathcal Z_s\) is empty.
Properness makes \(\pi(\mathcal Z)\) closed, so
\(U_S=S\setminus\pi(\mathcal Z)\) is an open neighborhood of \(s\)
whose fibers are all empty.  Thus the required neighborhood exists
in every case, without a flatness assumption.

Since \(S\) is integral, \(U_S\) contains its generic point, which has
characteristic zero.  To see explicitly that the open has
complex points, choose a nonempty principal open
\(D(h)\subset U_S\), with \(0\ne h\in R\).  The polynomial \(h\)
remains nonzero over \(\C\), so \(D(h)_\C\) is nonempty.  We may take
\(U=(U_S)_\C\).  The same invariance under field extension shows that
every complex fiber over \(U\) satisfies
\begin{equation}
 \dim V_{\PP^{v-1}_{\C}}(\nabla F_\theta)
 \le v-1-\delta.
 \label{eq:transfer-complex-gradient}
\end{equation}

Since \(\delta>0\), \(F_\theta\) cannot be the zero polynomial:
otherwise the left side of \eqref{eq:transfer-complex-gradient}
would be \(v-1\).  Over \(\C\), Euler's identity gives
\[
 rF_\theta=\sum_{i=1}^v x_i\partial_{x_i}F_\theta.
\]
Thus the ideals \((F_\theta,\nabla F_\theta)\) and
\((\nabla F_\theta)\) agree.  These equations describe the actual
projective hypersurface singular scheme.  Indeed, on the chart
\(x_i\ne0\), put \(z_j=x_j/x_i\) for \(j\ne i\) and
\(f_i=F_\theta/x_i^r\).  The hypersurface has equation \(f_i=0\)
and cotangent presentation
\[
 \cO_{V(f_i)}
 \xrightarrow{\,df_i\,}
 \cO_{V(f_i)}^{\,v-1}
 \longrightarrow\Omega_{V(f_i)/\C}\longrightarrow0.
\]
Its \((v-2)\)-nd Fitting ideal is generated by the local partial
derivatives of \(f_i\).  On \(f_i=0\), these are the homogeneous
gradient equations in this chart: Euler's identity supplies the
derivative in the omitted \(x_i\) direction.  The Jacobian criterion
therefore identifies this scheme with the singular scheme.
Equation~\eqref{eq:transfer-complex-gradient} proves
\eqref{eq:transfer-singular-bound}.
\end{proof}

All coefficients in \eqref{eq:transfer-projective-family} are integral;
in particular, the construction does not invert \(2\) or \(r\).
Division by \(r\) occurs only in characteristic zero.  In applying the
lemma to \(G_D\), we use its coefficient formula
\eqref{eq:selector-coefficient}, before multiplying by the projection
scalar \(a_0\), which vanishes in characteristic two.

\subsection{A smooth section of the required dimension}

Avoiding the singular locus of a hypersurface is one condition on a
linear section.  The section must also meet its smooth locus
transversely.  The following incidence calculation establishes both
conditions and keeps track of the projective dimension.

\begin{proposition}[Smooth linear section]\label{prop:smooth-slice}
Let \(F\) be a nonzero homogeneous complex polynomial of degree \(r\ge2\)
on a vector space \(V\) of dimension \(v\).  Set
\(Y=V_{\PP(V)}(F)\), and suppose
\[
 \dim\operatorname{Sing}Y\le v-1-\delta
\]
for an integer \(\delta\ge1\).  If \(2\le d<v\) and
\(d-1<\delta\), then a nonempty Zariski open subset of
\(\operatorname{Gr}(d,V)\) consists of subspaces \(H\) for which
\(F|_H\) is nonzero and defines a smooth degree-\(r\) hypersurface
in \(\PP(H)\).
\end{proposition}

\begin{proof}
Write \(\mathcal G=\operatorname{Gr}(d,V)\), so
\(\dim\mathcal G=d(v-d)\).  A point of \(\mathcal G\) is a vector
subspace of dimension \(d\), whose projectivization has dimension
\(d-1\).

First consider pairs \((p,H)\) with
\(p\in\operatorname{Sing}Y\) and \(p\in\PP(H)\).
For fixed \(p\), the choices of \(H\) form
\(\operatorname{Gr}(d-1,V/p)\), of dimension \((d-1)(v-d)\);
here \(p\) also denotes its line in \(V\).  This incidence has
dimension at most
\[
 (v-1-\delta)+(d-1)(v-d)
 =\dim\mathcal G+d-1-\delta
 <\dim\mathcal G.
\]
Its projection to \(\mathcal G\) is closed because
\(\operatorname{Sing}Y\) is projective.  It is a proper subset, so
outside it the section avoids every singular point of \(Y\).

For \(p=[x]\in Y_{\rm reg}\), the nonzero differential \(dF_x\)
has a kernel \(K_p\subset V\) of dimension \(v-1\).  Euler's
identity at \(F(x)=0\) shows that \(p\subset K_p\).
If \(p\subset H\), the restricted equation is singular at \(p\)
exactly when
\[
 p\subset H\subset K_p.
\]
The homogeneous partial derivatives define the Gauss morphism
\(Y_{\rm reg}\to\PP(V^*)\), \(p=[x]\mapsto[dF_x]\).
The pairs above form a closed incidence in \(Y_{\rm reg}\times\mathcal G\):
they are the inverse image of the closed flag incidence
\(p\subset H\subset\ker\ell\) under \((p,H)\mapsto(p,H,[dF_x])\).
For fixed \(p\), these subspaces form
\(\operatorname{Gr}(d-1,K_p/p)\), of dimension
\((d-1)(v-d-1)\).  Since \(\dim Y_{\rm reg}\le v-2\), this
tangency incidence has dimension at most
\[
 (v-2)+(d-1)(v-d-1)=d(v-d)-1.
\]
Its image in \(\mathcal G\) is constructible, and the closure of
that image has dimension at most \(\dim\mathcal G-1\).
Removing this closure and the preceding singular-point incidence
image leaves a nonempty open subset of the irreducible
Grassmannian.

For \(H\) in that open, every zero of \(F|_H\) has nonzero
differential on \(H\).  The restriction cannot vanish identically:
if it did, every point \(p\in\PP(H)\) would lie in \(Y\), singular
avoidance would put it in \(Y_{\rm reg}\), and differentiation of
the zero restriction would force \(H\subset K_p\), contrary to
the tangency exclusion.

Thus \(F|_H\) is a nonzero homogeneous polynomial of degree \(r\).
Since \(d\ge2\), it defines a nonempty hypersurface of dimension
\(d-2\) in \(\PP(H)\).  The nonzero differentials prove its
smoothness.
\end{proof}

\subsection{A smooth form obtained from the permanent}

We now apply the two geometric results to the coefficient polynomial.
The resulting form has linear degree and quadratically many variables,
and is an affine projection of a permanent of linear size.

\begin{samepage}
\begin{proposition}[A smooth form obtained from the permanent]
Let \(k\ge128\) be an integer\label{prop:smooth-projection}, and put
\[
 q=\left\lfloor\frac{k-1}{3}\right\rfloor,\qquad
 r=k-2q,\qquad d=\left\lfloor\frac{k^2}{100}\right\rfloor.
\]
There is a nonzero homogeneous polynomial
\(G\in\C[x_1,\ldots,x_d]\) of degree \(r\) such that its zero set in
\(\PP^{d-1}_{\C}\) is smooth and, for every integer \(m\ge11k-2q\),
there is an affine map \(\Phi_m:\C^d\to\Mat_m(\C)\) satisfying
\[
 G=\per_m\circ\Phi_m.
\]
The parameters also satisfy
\[
 r-1\ge\frac{k}{4},\qquad d-1\ge\frac{k^2}{200}.
\]
\end{proposition}
\end{samepage}

\begin{proof}
Set \(T=\binom{q+2}{2}\).  Since \(k\ge128\),
\[
 q\ge\frac{k-4}{3}\ge\frac{k}{4},\qquad
 T\ge\frac{q^2}{2}\ge\frac{k^2}{32}.
\]
Consequently
\begin{equation}
 \Delta:=T-k
 \ge\frac{k^2}{32}-k
 \ge\frac{k^2}{64}>0.
 \label{eq:transfer-codimension-size}
\end{equation}
Also \(T\le(k+2)(k+5)/18<k^2\) for these \(k\), so
\(1\le\Delta\le k^2\).

Choose pairwise distinct
\(\bar\theta_1,\ldots,\bar\theta_k\in\overline{\mathbb F}_2\)
and put \(\bar D=\operatorname{diag}(\bar\theta_1,\ldots,\bar\theta_k)\).
Proposition~\ref{prop:gradient-codimension} gives
\[
 \dim V_{\mathbb A^{k^2}_{\overline{\mathbb F}_2}}(\nabla G_{\bar D})
 \le k^2+k-T=k^2-\Delta.
\]
The coefficient formula \eqref{eq:selector-coefficient}, with the
\(\theta_i\) treated as indeterminates, is integral in these parameters
and homogeneous of degree \(r\) in the \(k^2\) matrix variables.
Moreover,
\[
 r-1=k-2q-1
 \ge\frac{k-1}{3}\ge\frac{k}{4},
\]
so \(r\ge2\).  Lemma~\ref{lem:characteristic-transfer} supplies a
nonempty open subset \(U\subset\mathbb A^k_{\C}\).  Choose
\(\theta=(\theta_1,\ldots,\theta_k)\in U(\C)\) and put
\(D=\operatorname{diag}(\theta_1,\ldots,\theta_k)\).  Then \(G_D\ne0\)
and
\[
 \dim\operatorname{Sing}V_{\PP^{k^2-1}_{\C}}(G_D)
 \le k^2-1-\Delta.
\]
The complex coefficient projection permits arbitrary diagonal entries
for this \(D\).

For \(d=\lfloor k^2/100\rfloor\), we have \(2\le d<k^2\), and
\eqref{eq:transfer-codimension-size} gives
\[
 d-1<\frac{k^2}{100}<\frac{k^2}{64}\le\Delta.
\]
Proposition~\ref{prop:smooth-slice} therefore supplies a subspace
\(H\subset\Mat_k(\C)\) of dimension \(d\) such that \(G_D|_H\)
is nonzero and defines a smooth degree-\(r\) hypersurface in \(\PP(H)\).
Choose a linear isomorphism \(\Lambda:\C^d\to H\) and set
\(G=G_D\circ\Lambda\).  This gives the asserted smooth form.  The
remaining scale estimate is
\[
 d-1\ge\frac{k^2}{100}-2\ge\frac{k^2}{200},
\]
where the last inequality holds for \(k\ge20\).

Finally, \(0\le3q\le k-1\), so \(0\le q\le\lfloor k/2\rfloor\).
For each integer
\(m\ge11k-2q\), Proposition~\ref{prop:selector} gives a matrix
\(A_m(X)\) of affine-linear forms with
\[
 \per_m(A_m(X))=a_0G_D(X),\qquad a_0\ne0.
\]
Compose \(A_m\) with \(\Lambda\) and multiply its first row by
\(a_0^{-1}\).  The resulting affine map
\(\Phi_m:\C^d\to\Mat_m(\C)\) satisfies
\(\per_m\circ\Phi_m=G\), as required.
\end{proof}
